
\medskip
\phantomsection\label{lem:sga3-VIB-10.9}%
\noindent\textbf{10.9.---}
\emph{In the situation recalled at the beginning of
\hyperref[par:sga3-VIB-10.0]{10.0}, let \(j\in I\), and let
\(G_j\) be an \(S_j\)-group of finite presentation. For
\[
  G^0=(G_j\times_{S_j}S)^0
\]
to be representable, it is necessary and sufficient that there
exist \(i\in I\), \(i\geq j\), such that
\[
  (G_i)^0=(G_j\times_{S_j}S_i)^0
\]
be representable.}

The condition is sufficient
\textup{(\hyperref[prop:sga3-VIB-3.3]{3.3})}.

Let us show that the condition is necessary. Indeed, \(G^0\) is
of finite presentation over \(S\) and open in \(G\)
\textup{(\hyperref[prop:sga3-VIB-3.9]{3.9})}.\footnote{Translator's note: by 3.9, \(G^0\) is open in \(G\) and
quasi-compact over \(S\); being open in \(G\), which is of finite presentation over \(S\), it is then of finite
presentation over \(S\). The Gille--Polo re-edition (2011), Expos\'e~VI\textsubscript{B}, makes this step
explicit, and obtains \(H_i\) below from \hyperref[lem:sga3-VIB-10.3]{10.3} and
\hyperref[par:sga3-VIB-10.1]{10.1}: the open immersion \(G^0\to G\) comes from an open immersion over some
\(S_i\), and the group structure of \(G^0\) from that of an open subscheme of \(G_i\) after increasing \(i\).}
By
\textup{(\hyperref[par:sga3-VIB-10.0]{10.0})}, there exist
\(i\in I\), \(i\geq j\), and an open group subscheme \(H_i\) of
\(G_i\) such that
\[
  H_i\times_{S_i}S=G^0.
\]
The structural morphism \(G^0\to S\) is connected \textup{(}that
is, has geometrically connected fibers
\textup{(\hyperref[prop:sga3-VIA-2.1.1]{VI\(_A\), 2.1.1}))}.
Then \textup{(}EGA~IV\(_3\), 9.3.3 and 9.7.7\textup{)}, after
increasing \(i\), we may suppose that the structural morphism
\(H_i\to S_i\) is connected,%
\begingroup
\renewcommand{\thefootnote}{\fnsymbol{footnote}}%
\footnote[1]{\label{note:sga3-VIB-10.9-Hi-base}
Translator's note: the French source prints \(H_i\to S\).
Since \(H_i\) is an \(S_i\)-group and the argument spreads out
the connectedness of \(G^0\to S\), the base here is \(S_i\).}%
\endgroup\ hence that the underlying space of \(H_i\) is precisely \((G_i)^0\)
\textup{(\hyperref[lem:sga3-VIB-3.10.1]{3.10.1})}, hence that \(H_i\) represents
\((G_i)^0\).

\medskip
\phantomsection\label{par:sga3-VIB-10.10}%
\noindent\textbf{10.10.---}
Recall two very useful special cases of the situation stated
in \hyperref[par:sga3-VIB-10.0]{10.0}:\footnote{Translator's note: see EGA~IV\(_3\), 8.1.2~a)
and c), to which the Gille--Polo re-edition (2011), Expos\'e~VI\textsubscript{B}, refers here.}

\smallskip
\phantomsection\label{par:sga3-VIB-10.10-a}%
\noindent\textup{a)}
Given a point \(x\) of a scheme \(X\), put
\(S_0=\operatorname{Spec}\mathbb Z\), and consider the
decreasing filtered family \((S_i)_{i\in I}\) of
affine open neighborhoods of \(x\). Then
\[
  S=\operatorname{Spec}\mathcal O_{X,x}.
\]
In particular, if \(x\) is the generic point of an integral
scheme \(X\), then
\[
  S=\operatorname{Spec}\kappa(x).
\]

\smallskip
\phantomsection\label{par:sga3-VIB-10.10-b}%
\noindent\textup{b)}
Put \(S_0=\operatorname{Spec}\mathbb Z\), and consider the
family \((\mathcal A_i)_{i\in I}\), preordered by inclusion,
of the finitely generated \(\mathbb Z\)-subalgebras of the
ring of an affine scheme \(S\). Since the
\(\mathcal A_i\) are Noetherian rings, this makes it possible
in many cases to pass from the Noetherian case to the general
case.

We shall now give two results concerning the special case
considered in \hyperref[par:sga3-VIB-10.10-a]{(a)}.%

\medskip
\phantomsection\label{prop:sga3-VIB-10.11}%
\noindent\textbf{Proposition 10.11.---}
\emph{Let \(S\) be an integral scheme with generic point
\(\eta\), \(G\) an \(S\)-group of finite presentation, \(X\) an
\(S\)-scheme of finite presentation, \(u\) and \(v\) two
\(S\)-morphisms from \(X\) to \(G\), and \(i:H\to G\) and
\(j:K\to G\) two monomorphisms of finite presentation. Then:}
\begin{enumerate}
\item[\textup{(i)}] \emph{there exists a nonempty open subset \(S'\) of \(S\) such that, if \(G'\)
  \textup{(}respectively \(H'\), \(X'\), \(u'\), \(v'\)\textup{)} denotes the restriction of \(G\)
  \textup{(}respectively \(H\), \(X\), \(u\), \(v\)\textup{)} over the open subset \(S'\) of \(S\), the
  functor \(\underline{\operatorname{Transp}}(u',v')\) \textup{(}respectively
  \(\underline{\operatorname{Centr}}(u')\), \(\underline{\operatorname{Centr}}_{G'}H'\)\textup{)} is
  representable by a closed \(S'\)-subscheme of finite presentation \(T'_1\) \textup{(}respectively
  \(C'_1\), \(C'_2\)\textup{)} of \(G'\);}
\item[\textup{(ii)}] \emph{suppose that \(K_\eta\) is closed in \(G_\eta\) \textup{(}which is the case if
  \(K\) is an \(S\)-group subscheme of \(G\) \textup{(\hyperref[cor:sga3-VIB-1.4.2]{1.4.2}))}; then there
  exists a nonempty open subset \(S'\) of \(S\) such that, if \(G'\) \textup{(}respectively \(H'\),
  \(K'\)\textup{)} denotes the restriction of \(G\) \textup{(}respectively \(H\), \(K\)\textup{)} over the
  open subset \(S'\) of \(S\), the functor \(\underline{\operatorname{Transp}}_{G'}(H',K')\)
  \textup{(}respectively \(\underline{\operatorname{Norm}}_{G'}K'\)\textup{)} is representable by a closed
  \(S'\)-subscheme of finite presentation \(T'_2\) \textup{(}respectively \(N'\)\textup{)} of \(G'\);}
\item[\textup{(iii)}] \emph{suppose that \(H_\eta\) and \(K_\eta\) are closed in \(G_\eta\) \textup{(}which
  is the case if \(H\) and \(K\) are group subschemes of \(G\)
  \textup{(\hyperref[cor:sga3-VIB-1.4.2]{1.4.2}))}; then there exists a nonempty open subset \(S'\) of
  \(S\) such that, if \(G'\) \textup{(}respectively \(H'\), \(K'\)\textup{)} denotes the restriction of
  \(G\) \textup{(}respectively \(H\), \(K\)\textup{)} over the open subset \(S'\) of \(S\), the functor
  \(\underline{\operatorname{Transp}}^{\mathrm{str}}_{G'}(H',K')\) is representable by a closed
  \(S'\)-subscheme of finite presentation of \(G'\).}
\end{enumerate}

\emph{Proof.}
Since \(G_\eta\), \(X_\eta\), \(H_\eta\), \(K_\eta\) are flat over the field \(\kappa(\eta)\), and since
\(G_\eta\) is separated over \(\kappa(\eta)\)
\textup{(\hyperref[sga3:VIA:0.3]{VI\(_A\), 0.3})},\footnote{Translator's note: the 1970 text again refers to
VI\(_A\), 0.2 (see the note to \hyperref[prop:sga3-VIB-5.2]{5.2}); and the flatness over \(S'\) obtained below
comes from EGA~IV\(_3\), 11.2.6 (see \hyperref[par:sga3-VIB-10.1]{10.1}). The Gille--Polo re-edition (2011),
Expos\'e~VI\textsubscript{B}, cites VI\(_A\), 0.3 and 10.1.} by \textup{(}\hyperref[lem:sga3-VIB-10.2]{10.2},
\hyperref[par:sga3-VIB-10.10-a]{10.10~a)}, EGA~IV\(_3\), 8.8.2 and 8.10.5\textup{)} there exists an affine open
subset \(S'\) of \(S\) such that \(G'\) is an \(S'\)-group of finite presentation, flat and separated over
\(S'\), and such that \(X'\), \(H'\) and \(K'\) are flat and of finite presentation over \(S'\). If
\(H_\eta\) \textup{(}respectively \(K_\eta\)\textup{)} is assumed closed in \(G_\eta\), one can choose \(S'\) so
that \(H'\) \textup{(}respectively \(K'\)\textup{)} is closed in \(G'\). By EGA~IV\(_3\), 11.3.15, after
replacing \(S'\) by an affine open subset of \(S'\), one may suppose that \(G'\), \(H'\), \(K'\) and \(X'\) are
essentially free over \(S'\) \textup{(}in the sense of VIII,
\SGARefLink{sga3:viii:definition:6:1}{6.1}\textup{)}.\footnote{Translator's note: let \(A\) be the ring of
\(S'\), an integral domain. Each of these schemes, being of finite presentation over \(S'\), is covered by
finitely many affine open subschemes, whose rings \(B\) are flat \(A\)-algebras of finite presentation. For each
\(B\), EGA~IV\(_3\), 11.3.15 gives \(f_1,\ldots,f_n\in A\), generating the unit ideal, such that
\(B_{f_m}/(f_1,\ldots,f_{m-1})B_{f_m}\) is free over \(A_{f_m}/(f_1,\ldots,f_{m-1})A_{f_m}\) for every \(m\); if
\(f_m\) is the first nonzero \(f_k\), then \(B_{f_m}\) is free over \(A_{f_m}\). One then replaces \(S'\) by
\(D(f)\), \(f\) being the product of the elements so obtained, which is nonzero. The Gille--Polo re-edition
(2011), Expos\'e~VI\textsubscript{B}, details the same step over a noetherian ring, with EGA~IV\(_2\), 6.9.2 (see
the Addition below).} It then follows from \textup{(}VIII, \SGARefLink{sga3:viii:examples:6:5:item:b:01}{6.5~(b)}
and \SGARefLink{sga3:viii:examples:6:5:item:e:01}{(e)}\textup{)}%
\begingroup
\renewcommand{\thefootnote}{\ensuremath{(\ast)}}%
\footnote[1]{\label{note:sga3-VIB-10.11-star}See the note \((\ast)\) to the proof of
\hyperref[note:sga3-VIB-6.2-star]{6.2}; the 1970 text refers to it as the note on page~40, the page of this Expos\'e
on which it stands. Translator's note: the 1970 text cites ``VIII 5.5 b) et e)'' here, a misprint for VIII, 6.5~(b)
and (e), which the proof of 6.2 uses in the same way.}%
\endgroup{} that, under the hypotheses of the statement, the functors considered are representable by closed
subschemes \textup{(}hence of finite presentation over \(S'\)\textup{)} of \(G'\).

\par\medskip\noindent
\emph{Addition} (after the Gille--Polo re-edition (2011), Expos\'e~VI\textsubscript{B}; our formulation).
(a) When \(S'\) is not noetherian, a closed subscheme of \(G'\) need not be of finite presentation over \(S'\);
those obtained above are, for the following reason. In the proof of VIII,
\SGARefLink{sga3:viii:theorem:6:4}{6.4}, when the base is affine with ring \(A\) and \(Z\) is affine with ring
\(B\) free over \(A\), the closed subscheme representing \(\prod_{Z/S}Y/Z\) is defined by the ideal of \(A\)
generated by the coordinates of a family of generators of the ideal of \(Y\) in \(B\); so it is of finite
presentation when that ideal is of finite type, and the same holds when \(Z\) is covered by finitely many such
affine open subschemes. In the proof above, the schemes to which VIII, 6.4 is applied through VIII, 6.5 are of
finite presentation over \(G'\), hence quasi-compact over it, and the closed subschemes in question (inverse
images of the diagonal of \(G'\), or of \(K'\), or of \(H'\)) are of finite presentation, \(G'\) being
separated and of finite presentation over \(S'\), and \(H'\), \(K'\) of finite presentation. The re-edition
obtains this finiteness differently: by \hyperref[par:sga3-VIB-10.10-b]{10.10~b)} it descends \(G'\), \(X'\),
\(H'\), \(K'\) and the morphisms between them to the spectrum of a noetherian subring \(A_0\) of the ring of
\(S'\), proves the result there, using the generic freeness theorem EGA~IV\(_2\), 6.9.2 (valid over a noetherian
integral ring) in place of 11.3.15, and comes back to \(S'\) by base change, the functors considered being
compatible with base change.
(b) The re-edition states the proposition for morphisms of \(S\)-schemes \(u,v,i\colon Y\to G\) and
\(j\colon Z\to G\) of which only the generic fibres \(i_\eta\) and \(j_\eta\) are assumed to be closed
immersions: \(i\) and \(j\) are then closed immersions over a nonempty open subset of \(S\)
(\hyperref[par:sga3-VIB-10.1]{10.1}), and the statement above applies. It treats separately, as its
Corollary~10.11.1, two quasi-compact monomorphisms of \(S\)-groups \(H\to G\) and \(K\to G\), with \(H\) and \(K\)
of finite presentation: their generic fibres are closed immersions
(\hyperref[cor:sga3-VIB-1.4.2]{1.4.2}), and the centralizer and normalizer of \(H'\) are then represented by
closed subgroup schemes of \(G'\).
